\chapter*{Preface}
\addcontentsline{toc}{chapter}{Preface}
\markboth{Preface}{Preface}

Here is the moment this book was born.

Somebody asked me why their agent took forty seconds to do four things that each
took a tenth of a second, and I could not tell them. Not because the answer was
subtle, but because nobody in the room had any instrument that could see where
the time went. We had a stopwatch and a shrug.

I have since watched that same conversation happen at four companies. The
details differ. The shape never does: a system that works, a number nobody can
explain, and a team optimising the one layer their existing monitoring happens
to cover. Chapter~\ref{ch:primer} tells the full story and then takes the forty
seconds apart.

This book is the instrument I wanted that afternoon.

\section*{What this book is}

This is a protocol book that is secretly a performance book. Or possibly a
performance book that is secretly a protocol book. The two turn out to be the
same subject, and that is the whole thesis.

The model here is Ilya Grigorik's \emph{High Performance Browser Networking},
which did something I have never stopped admiring. It taught web developers TCP.
Not trivia about TCP. Actual TCP: slow start, congestion windows, the three-way
handshake, why your first byte arrives when it arrives. And it justified every
page of that by cashing it out into things you could do on Monday morning.
Nobody had to be persuaded that the protocol layer mattered, because every
chapter ended with a measurement that only made sense if you understood the
layer underneath.

I want to do that for the Model Context Protocol. Not because MCP is beautiful,
though parts of it are, but because you cannot build a good agentic application
without knowing what the protocol does on the wire. The abstractions leak
immediately and constantly. A tool description is not free; you pay for it on
every single request, forever. An extra round trip is not a small cost; it is
usually the largest single item on the bill. A cache hint you ignored is a
hundred milliseconds you spent for nothing. None of this is visible from
inside a framework, and all of it is obvious once you have watched the bytes.

So: three parts physics, three parts protocol mechanics, and a lot of numbers.

\section*{Who it is for}

You have shipped something. Maybe it is a chat assistant with tools bolted on,
maybe it is an autonomous pipeline that runs at three in the morning, maybe it
is a server that other people's agents call. It works, mostly. You have a
nagging sense that you are holding it wrong.

You do not need to know MCP already. Chapter~\ref{ch:architecture} builds it up
from nothing. You do need to be comfortable reading code, thinking about
latency, and being told that a thing you are currently doing is expensive.

If you are looking for a tutorial that gets you to a working server in ten
minutes, this is not it, and the official documentation does that job well.
This book is for the part afterwards, when it works and you need it to be
\emph{good}.

\section*{The version problem, and what I did about it}

Every book about a young protocol has the same terrifying paragraph in the
preface, where the author admits the thing might change and asks for your
patience. Here is mine, except I think the situation is better than usual, and
for a specific reason.

This book is pinned to revision \texttt{2026-07-28}. That revision is the
largest change since MCP launched. It deleted the initialization handshake. It
deleted protocol-level sessions. It deleted the standalone SSE endpoint,
stream resumability, and server-initiated requests. It deprecated Sampling,
Roots, and Logging. If you learned MCP in 2025, a meaningful fraction of what
you know is now wrong, and I will point out exactly which fraction as we go.

But the same revision also adopted a formal feature lifecycle: Active,
Deprecated, Removed, with a minimum twelve-month window between the middle one
and the last one. That is the first version of MCP that comes with a promise about how it will
break. MCP will break. The promise is that it will break \emph{slowly and in
public}, which is the only kind of promise a protocol can honestly make and the
only kind that lets you plan.

So the deprecated features are in this book, boxed and labelled, because you
have to maintain systems that use them. They are never presented as something
to adopt. When you see one of these:

\begin{legacybox}{Sampling}
The 2025 design let a server ask the client to run a completion on its behalf.
It is Deprecated as of \texttt{2026-07-28} and eligible for removal no earlier
than 2027-07-28. Chapter~\ref{ch:mrtr} covers what it was for, why it lost, and
what to do with the servers that still speak it.
\end{legacybox}

\noindent you are reading about the past, on purpose, so you can recognise it in
the wild.

\section*{Meridian, and why the numbers are real}

Technical books have a bad habit. They say things like ``this reduces latency
significantly'' and then move on, and you are left with a claim you cannot
check and a number you cannot use.

I did not want to write that book, so this one comes with a program.

Meridian is a financial analytics assistant: four MCP servers, one host, an
agent loop, a load-test harness, and a measurement rig. It ships in the
repository. It has no third-party dependencies, which means it also happens to
be a complete implementation of \texttt{2026-07-28} in about 3,900 lines of
Python, and every listing in this book is lifted from it. When
Chapter~\ref{ch:tools} claims a bloated tool catalogue costs you 4,371 extra
tokens on every model turn, that number came out of
\texttt{meridian/bench/results.json}, and you can regenerate it with
\texttt{make bench} and watch it change when you edit the catalogue.

Two honest caveats, stated here and repeated wherever they matter.

First, model inference is simulated. A book cannot ship a reproducible large
language model, so Meridian uses a stub with a stated latency distribution.
Everything downstream of inference (serialisation, transport, server execution,
caching, round-trip counts) is measured for real.

Second, the transport numbers are loopback numbers. Localhost has no physics
worth speaking of. That is deliberate: it isolates protocol overhead so you can
see it, and then Chapter~\ref{ch:measuring} shows you how to add your own
network back in. A book that mixed the two would be measuring my office
wi-fi and calling it your architecture.

\begin{measurebox}{What a box like this means}
Numbers in boxes like this one came from a real run of the harness on the
machine described in Chapter~\ref{ch:primer}. If a number is modelled rather
than measured, the box says so explicitly.
\end{measurebox}

\section*{A note on how this is written}

Long. Opinionated. Occasionally funny, when the joke earns its place.

I have tried very hard to avoid the register that technical writing drifts into
when nobody is watching, the one where every paragraph restates the previous
paragraph in slightly different words and nothing is ever asserted plainly. If
I think something is a bad idea, I say it is a bad idea and then show you the
measurement. If I am guessing, I say I am guessing.

Each chapter opens with a line of film dialogue. This is not decoration. It is
a compression trick: the right quote does about a paragraph of work in eleven
words, and it tells you what the chapter is actually arguing before the
argument starts.

Also, there is not a single em dash in this book. This started as a constraint
somebody imposed on me and became a genuine discipline. Em dashes let you
staple two thoughts together without deciding how they relate. Take them away
and you have to pick: is this a list, a consequence, an aside, or two
sentences? Usually it was two sentences. The prose got better. Make of that
what you will.

\section*{Thank you}

To the MCP maintainers and the SEP authors, whose design discussions are public
and unusually candid, and who deprecated three of their own features in one
release rather than defend them. That is rarer than it should be.

To everybody who has shipped an agent that was mysteriously slow and then let
me look at the traces.

\vspace{14pt}
\noindent\emph{Krimler}\\
\emph{July 2026}

\cleardoublepage
